\documentclass{situationist-preprint}

\title{Grounding Without a Fundamental Level}
\msid{MET-10-a}
\msfamily{MET-10}
\msdiscipline{Metaphysics}
\mstype{Principal manuscript}
\mskeywords{fundamentality; grounding; levels; infinite regress; metaphysical infinitism; gunk; well-foundedness}

\begin{document}
\maketitle

\begin{abstract}
Is there a fundamental level of reality, a level of objects whose features are grounded in nothing further and in which everything else is grounded? We argue that there is not. Every level of reality is grounded in a lower one, and the resulting infinite descent is not vicious. The argument has two parts. First, given that the facts about a composite at any level are grounded in facts about its parts and their relations, and that every level of parts has a finer level, no level is fundamental; we defend the divisibility premise from the continuity of physical fields and the record of physics. Second, we distinguish three ways in which a regress can be vicious, by inconsistency, by leaving some member without a full ground, and by answering a general question with an instance of what it asks about, and show that the descent of levels has none of them. The descent concerns objects and the grounding of wholes in their parts; it is compatible with a fixed stock of genuine properties instantiated at every level, with grounding as a strict partial order that is not well-founded, and with the claim of MET-09 that the cosmos, which is not a level, is the one fundamental object.
\end{abstract}

\section{The question}\label{s:question}

The world presents itself in levels. Organisms are made of cells, cells of molecules, molecules of atoms, atoms of nucleons and electrons, nucleons of quarks. At each step, the features of the things at one level are explained by the features of their parts and the relations among them: the rigidity of a table by the bonds between its molecules, the chemistry of an element by the structure of its atoms. The hierarchy was made into an explicit picture of the sciences by \citet{OppenheimPutnam1958}, for whom each level is a decomposition of the level above it into parts, and the unity of science consists in the reduction of each level to the one below.

The question is whether the hierarchy ends. \emph{Metaphysical foundationalism} holds that it does: there is a level of objects, the fundamental level, whose features are not grounded in anything, and every other fact is grounded, ultimately, in facts about that level. The atomists held a version of this view, and so did Leibniz and the logical atomists. \emph{Metaphysical infinitism} holds that it does not: every level is grounded in a lower one, and the descent goes on without end \citep[see][]{Morganti2014}. A third position holds that the question has no answer we can know; \citet{Schaffer2003} argued that the evidence of physics, which has repeatedly found a further level beneath the one it took to be the last, gives us no reason to believe in a fundamental level, and treated infinite descent as an open epistemic possibility.

We defend infinitism.

\begin{quote}
\textbf{Thesis.} There is no fundamental level of reality. Every level is grounded in a lower one, and this infinite descent is not vicious.
\end{quote}

The thesis has two parts, and the second is the more contested. Many philosophers grant that an infinite descent of parts, or \emph{gunk}, is possible \citep{Lewis1991,Zimmerman1996}, but hold that an infinite descent of \emph{grounds} is not, because, as \citet{Schaffer2010} put the worry, being would then be infinitely deferred and never achieved. We argue that this worry rests on a picture of grounding as the transmission of a quantity, and that on the correct picture, on which grounding is a relation of metaphysical explanation among facts each of which obtains, the descent leaves nothing unexplained that a foundationalist explains.

Two clarifications fix the scope of the thesis. First, the descent concerns \emph{objects} and the grounding of facts about wholes in facts about their parts. It does not require a descent of \emph{properties}, in which every property would be grounded in a more basic one; the same genuine properties, in the sense of MET-02-a, may be instantiated at every level. Second, a level, in the sense of the thesis, is a level of decomposition of the world into parts. The world as a whole is not one of its own levels, and the thesis is therefore silent on whether the whole is fundamental. MET-09 argues that it is, and Section~\ref{s:cosmos} shows that the two theses fit together.

Section~\ref{s:history} reviews the history of the question. Section~\ref{s:framework} sets out the framework, Section~\ref{s:argument} the argument, and Section~\ref{s:vicious} the account of viciousness. Section~\ref{s:cosmos} relates the result to MET-09, Section~\ref{s:cases} considers further cases, and Section~\ref{s:objections} answers objections.

\section{A history of the question}\label{s:history}

\paragraph{Aristotle.} The classical source of the demand for a first term is \emph{Metaphysics} $\alpha$~2, where Aristotle argues that the causes of things are not infinite either in series or in kind: in a series of intermediates, the intermediates are causes only because there is a first, and if there were no first there would be no cause at all. The argument concerns efficient, material, final and formal causes alike, and its premise is that an intermediate is a cause only derivatively, by transmitting the causality of what is before it. The case against the argument is the series of intermediates that are causes in their own right. In a chain of grounds in which each member fully grounds the next, each member is a ground whether or not there is a first; Aristotle's premise holds for instrumental causes, such as the stick moved by the hand, but not for grounds in general.

\paragraph{Leibniz.} Leibniz argued that composites presuppose simples: there must be simple substances, because there are composites, and a composite is nothing but a collection of simples (\emph{Monadology} \S\S\,1--3). In his correspondence with Arnauld he put the point in terms of unity: what is not truly one being is not truly a being either, so that if every body were divisible without end, there would be aggregates of aggregates and nothing of which they were aggregates. The defeating case is the continuum. A line segment is a whole whose parts are segments, every one of which has segments as parts; the segment is not thereby a collection of nothing. Leibniz's premise, that the being of an aggregate is borrowed from the being of its ultimate members, assumes that there are ultimate members to lend it, which is the conclusion.

\paragraph{Kant.} The second antinomy sets the thesis that every composite substance consists of simple parts against the antithesis that no composite thing consists of simple parts (\emph{Critique of Pure Reason} A434/B462). Kant took both to be provable, and concluded that the question cannot be answered of things in themselves. His proof of the thesis repeats Leibniz's: if composition were removed in thought, nothing would remain unless there were simples. The proof of the antithesis rests on the infinite divisibility of space. The lesson we draw is not Kant's. The thesis proof assumes that removing composition in thought must leave a residue, and an infinitist denies that; the antithesis proof is an argument from the structure of space, which we take up in Section~\ref{s:argument}.

\paragraph{Logical atomism.} Russell's logical atomism held that the analysis of complexes terminates in simples, and that the world consists of atomic facts about them \citep{Russell1918}. The commitment was largely methodological: analysis must stop somewhere if it is to be completed. Russell himself allowed that there might be no simples, and that the theory might hold only of the units at which analysis in fact stops. The defeating case was the development of physics after 1918, in which successive candidates for the smallest units turned out to have structure.

\paragraph{Levels and exclusion.} In the philosophy of mind the hierarchy of levels acquired a new role. \citet{Kim1998} argued that if mental properties are realized by physical ones, the causal powers of the mental are excluded by those of their physical realizers, which are sufficient for the same effects. \citet{Block2003} pointed out that the argument generalizes: if the powers of each level are excluded by those of the level beneath, and there is no bottom level, causal powers drain away and nothing has any. The drainage argument is the most important recent argument that the descent, if there is one, would be vicious, and we answer it in Section~\ref{s:cases}.

\paragraph{The recent debate.} \citet{Schaffer2003} argued that there is no good evidence for a fundamental level, and that the history of physics supports at most the conclusion that each level so far has turned out to have a further level beneath it. \citet{Cameron2008} examined the arguments that infinite descent of grounds is vicious and found none of them compelling, but argued that a well-founded structure is preferable because it provides a unified explanation of all that is derivative; on his view, well-foundedness is a contingent truth we have theoretical reason to believe. \citet{Bliss2013} distinguished the senses in which a regress might be vicious and argued that it is much harder than commonly supposed to show that a descending chain of grounds is vicious in any of them. \citet{Morganti2014} defended metaphysical infinitism, \citet{Tahko2014} examined which kinds of infinite descent are compatible with an ordinary conception of fundamentality, and \citet{Dixon2016} showed that \enquote{well-founded} admits several inequivalent readings, of which the requirement that there be no infinite descending chains of grounds is only one. We side with Bliss and Morganti on viciousness, answer Cameron's argument from unified explanation in Section~\ref{s:objections}, and use Dixon's distinction in Section~\ref{s:cosmos}.

\section{Framework}\label{s:framework}

We speak of facts, of objects, and of the relation of partial grounding between facts. We write $f \prec g$ for \enquote{$f$ partly grounds $g$}. A non-empty set of facts $\Gamma$ \emph{fully grounds} $g$ if $g$ obtains in virtue of the members of $\Gamma$ taken together, with nothing further needed; each member of a full ground of $g$ partly grounds $g$. We do not consider zero-grounding, the limiting case of a fact grounded in the empty set. By the \emph{cosmos} we mean the fusion of all concrete objects.

\begin{definition}[Level]\label{d:level}
A \emph{level} is a plurality of concrete objects, each of them a proper part of the cosmos, that do not overlap and that together compose the cosmos. A level $L'$ is \emph{finer} than a level $L$ if every member of $L'$ is a part of some member of $L$ and some member of $L$ has at least two members of $L'$ as parts.
\end{definition}

\begin{definition}[Level facts]\label{d:facts}
The \emph{facts at} a level $L$ are the facts about the members of $L$: which objects they are, which genuine properties they have, which genuine relations they stand in to one another, and which pluralities of them compose which objects.
\end{definition}

\begin{definition}[Grounded level; fundamental level]\label{d:grounded}
A level $L$ is \emph{grounded in} a level $L'$ if every fact at $L$ has a full ground consisting of facts at $L'$. A level is \emph{fundamental} if no fact at it is partly grounded in any fact at any other level.
\end{definition}

Definition~\ref{d:level} makes a level a decomposition of the cosmos into proper parts. Two consequences matter. The cosmos itself is not a level, since it is not a proper part of itself; and levels need not correspond to the levels of the special sciences, since a level may contain a quark in one region and a galaxy in another. The second consequence is harmless. The hierarchy of the sciences is a selection of levels that are uniform in scale, and our results hold for every level and so for those.

\begin{assumption}[Grounding]\label{a:ground}
Partial grounding is a strict partial order on facts: no fact partly grounds itself, and if $f \prec g$ and $g \prec h$, then $f \prec h$.
\end{assumption}

Assumption~\ref{a:ground} is the assumption of MET-12-a, and it is standard \citep{Fine2012,Rosen2010,Schaffer2009}. As MET-12-a notes, it does not require that grounding be well-founded: it is consistent with there being an infinite sequence of facts each of which is partly grounded in the next. We assume, with MET-14, that grounding is a single relation, whose apparent varieties are that relation applied at different levels; nothing below depends on whether the grounding of a whole in its parts is the same relation as, for instance, the grounding of a disjunction in a disjunct, but the unity of grounding makes the transitivity clause of Assumption~\ref{a:ground} unrestricted.

\begin{principle}[Aggregation]\label{p:aggregation}
If $L'$ is finer than $L$, every fact at $L$ has a full ground consisting of facts at $L'$.
\end{principle}

\begin{principle}[Divisibility]\label{p:divisibility}
For every level there is a finer level.
\end{principle}

\begin{principle}[Fixed stock]\label{p:stock}
There is a set of genuine properties and relations such that the genuine properties and relations instantiated at every level belong to it.
\end{principle}

\ref{p:aggregation} says that a whole is no more than its parts, their genuine properties and relations, and the fact that they compose it. It includes relations: the facts at the finer level include the spatiotemporal relations among the parts and, if physics requires them, the non-separable relations of quantum entanglement, as well as the composition facts that MET-06 states in terms of common fate. Since the facts at $L'$ include which pluralities of members of $L'$ compose which objects, they include the existence of the members of $L$. The principle is the metaphysical core of the practice of explaining wholes by their parts, and it is accepted by foundationalists and infinitists alike. It is denied by those who hold that some wholes have strongly emergent features, not grounded in their parts and their relations; we consider them in Section~\ref{s:objections}.

\ref{p:divisibility} is the substantive premise, and Section~\ref{s:divisibility} defends it. \ref{p:stock} is not needed for the argument that there is no fundamental level. We state it because it records what the descent does \emph{not} require, and because, as Section~\ref{s:vicious} shows, it helps to see why the descent is not vicious. It is the scientific conception of sparse properties of \citet{Schaffer2004}, on which the genuine properties are those that physics finds at every scale, such as mass, charge and spatiotemporal separation, as opposed to the fundamental conception, on which they are the properties of the bottom level.

\section{The argument}\label{s:argument}

\begin{lemma}\label{l:finer}
If $L'$ is finer than $L$, then $L$ is grounded in $L'$.
\end{lemma}

\begin{proof}
Immediate from \ref{p:aggregation} and Definition~\ref{d:grounded}.
\end{proof}

\begin{lemma}\label{l:notfund}
No level that has a finer level is fundamental.
\end{lemma}

\begin{proof}
Let $L'$ be finer than $L$. By Definition~\ref{d:level}, some member $m$ of $L$ has at least two members of $L'$ as parts, so $L' \neq L$. The fact that $m$ exists is a fact at $L$. By Lemma~\ref{l:finer} it has a full ground consisting of facts at $L'$, and since full grounds are non-empty, some fact at $L'$ partly grounds it. So a fact at $L$ is partly grounded in a fact at another level, and $L$ is not fundamental.
\end{proof}

\begin{theorem}[No fundamental level]\label{t:nofund}
Every level is grounded in a lower one, and no level is fundamental.
\end{theorem}

\begin{proof}
Let $L$ be a level. By \ref{p:divisibility} there is a level $L'$ finer than $L$; by Lemma~\ref{l:finer}, $L$ is grounded in $L'$; by Lemma~\ref{l:notfund}, $L$ is not fundamental.
\end{proof}

\begin{corollary}[Descent]\label{c:descent}
There is an infinite sequence of levels $L_0, L_1, L_2, \dots$, each finer than and grounding the one before it, and an infinite sequence of facts $f_0, f_1, f_2, \dots$ with $f_{n+1} \prec f_n$ for every $n$. Partial grounding is therefore not well-founded.
\end{corollary}

\begin{proof}
Choose any level $L_0$ and, by \ref{p:divisibility}, a finer level $L_{n+1}$ for each $L_n$. Choose any fact $f_0$ at $L_0$ and, for each $n$, a member $f_{n+1}$ of a full ground of $f_n$ consisting of facts at $L_{n+1}$, which exists by Lemma~\ref{l:finer}. Then $f_{n+1} \prec f_n$. By transitivity the facts are pairwise distinct, since $f_m \prec f_n$ for $m > n$ and grounding is irreflexive, so the sequence is infinite and has no least member.
\end{proof}

The choice of a finer level at each step uses the axiom of dependent choice, which we take for granted.

\begin{proposition}[Fixed stock is consistent with descent]\label{p:stockcons}
\ref{p:aggregation}, \ref{p:divisibility} and \ref{p:stock} are jointly consistent, together with Assumption~\ref{a:ground}.
\end{proposition}

\begin{proof}
Consider a model in which the cosmos is a cube of homogeneous matter of uniform mass density, the levels are the decompositions of the cube into congruent cubes of side $2^{-n}$, and the genuine properties are mass, given by volume times density, and the spatial relations among cubes. Let the facts at level $n$ be the masses of the cubes at that level, their spatial relations and the composition facts, and let a fact at level $n$ be partly grounded in exactly the facts at levels $m > n$ about cubes that are parts of the cubes it concerns. This relation is irreflexive and transitive. Every fact at level $n$ is fully grounded by the facts about the eight sub-cubes of each cube it concerns at level $n+1$, since the mass of a cube is the sum of the masses of its eight sub-cubes and its spatial relations are fixed by theirs; so \ref{p:aggregation} holds. Every level has a finer one, so \ref{p:divisibility} holds. At every level the only genuine properties instantiated are mass and the spatial relations, so \ref{p:stock} holds.
\end{proof}

The model is a toy, but it displays the point recorded in MET-02-a: an infinite descent of objects is compatible with a fixed stock of genuine properties, instantiated at every level. What the sparse theory of universals requires is that every property be constructed from genuine properties, and that requirement is met at every level; it does not require that grounding terminate in a bottom level of objects.

\subsection{Defence of divisibility}\label{s:divisibility}

\ref{p:divisibility} is an empirical claim about the actual world, and the case for it is empirical. We give three considerations.

\emph{Continuity.} Our best physical theories describe matter by fields defined on a continuous spacetime. The portion of a field in a region has as parts its portions in the subregions, and every region has proper subregions. If the matter in a region is the field in that region, every material object has proper parts. The point-sized portions are not a bottom level, for two reasons. The value of a field at a point is not an intrinsic feature of anything located at the point: the physically significant quantities, such as gradients and velocities, depend on the field in every neighbourhood of the point, a point made against what he called pointillisme by \citet{Butterfield2006}. And in quantum field theory the field operators are not defined at points at all; they are operator-valued distributions, which yield operators only when smeared over regions \citep{StreaterWightman1964}. There is no level of physical facts at points for the descent to terminate in.

\emph{The record.} \citet{Schaffer2003} pointed out that every level that physics has proposed as the last has turned out to have structure: atoms, nuclei, nucleons. The inference from this record to infinite descent is inductive and defeasible, and we do not rest the case on it alone. But the record shifts the burden: a foundationalist who holds that the current level is the last makes a claim that every previous instance of it has falsified.

\emph{Point particles.} The standard model treats electrons and quarks as point particles without internal structure. This is not evidence of a bottom level. The particles of the standard model are excitations of fields, and the claim that they are pointlike is the claim that no internal structure has been found at the energies so far probed, which is consistent with continuity of the fields of which they are excitations. An electron is not an object with no proper parts; it is a state of the electron field, which, like every field, has its portions in every subregion of the region over which the state is spread.

We conclude that \ref{p:divisibility} is true, and that the world has no fundamental level. The conclusion is a posteriori and contingent: we do not claim that a world with a fundamental level is impossible. MET-06-a states its criterion of composition in terms of simples but notes that, if matter is gunky, the criterion can be relativized to any level of decomposition and applied at each level; on the present result, that is the relativization that applies to the actual world.

\section{Why the descent is not vicious}\label{s:vicious}

A regress may be harmless or harmful. We say that an infinite sequence of facts is \emph{descending} if each member is partly grounded in the next. In the terminology of MET-22-a (Definition~2.2), a sequence in which each fact obtains in virtue of the next is called \emph{vicious}; MET-22-a needed the term only to contrast Bradley's regress, which it showed to run in the opposite direction, and it left open whether descending chains of grounds are harmful in general (Section~3). To avoid a merely verbal disagreement, we call such sequences \emph{descending} and reserve \emph{vicious} for those that are harmful. The question is what harm a descending sequence could do.

\begin{definition}[Vicious regress]\label{d:vicious}
A descending sequence of facts is \emph{vicious} if it has one of the following features:
\begin{enumerate}[label=(V\arabic*)]
\item \emph{inconsistency}: the principles that generate it, together with Assumption~\ref{a:ground}, are inconsistent;
\item \emph{groundlessness}: some member that is not fundamental has no full ground;
\item \emph{self-presupposition}: it is generated by answering a general question about a kind of fact, of the form \enquote{how does any fact of kind $K$ obtain?}, with a fact of kind $K$, so that the answer presupposes what was asked.
\end{enumerate}
\end{definition}

The three features are the ones the literature has identified \citep{Nolan2001,Bliss2013}. (V1) is the strongest: a regress that cannot exist cannot be actual. (V2) is the failure that the foundationalist fears: that in the descent the obtaining of each member is never accounted for, because each waits on another that is never secured. (V3) is the classical form of a vicious explanatory regress. If one explains how a relation relates its terms by citing a further relation that relates it to them, the explanation cites an instance of what was to be explained, and the regress shows that the explanation fails at every step \citep{Nolan2001}; MET-22-a shows that this is the reading of Bradley's regress on which it is vicious, and that the premise generating it is false.

\begin{theorem}[Non-viciousness]\label{t:benign}
The descent of Corollary~\ref{c:descent} is not vicious.
\end{theorem}

\begin{proof}
(V1) By Proposition~\ref{p:stockcons}, the principles generating the descent are consistent with Assumption~\ref{a:ground}.

(V2) Every member $f_n$ of the descent is a fact at a level $L_n$, and by Lemma~\ref{l:finer} it has a full ground consisting of facts at $L_{n+1}$. So no member lacks a full ground.

(V3) The descent is generated by \ref{p:aggregation} and \ref{p:divisibility}, which answer particular questions: why this object has this mass, why these objects stand in this relation. They answer each by citing particular facts about other objects, the parts. Neither principle answers a general question of the form \enquote{how does any fact of kind $K$ obtain?}. In particular, the descent does not purport to explain why anything has mass at all, or why anything exists at all, by citing a further fact of the same kind. The general question why there is something rather than nothing is not answered by the descent and is not meant to be; it is the question of MET-01.
\end{proof}

The proof of (V2) is where the foundationalist will press. Each member has a full ground, she will say, but the full ground has a full ground, and so on, and the obtaining of the first member is thereby made to depend on infinitely many facts, none of which obtains on its own. The objection has a true premise and a false conclusion. Each member does depend on infinitely many facts, and none obtains ungroundedly. But every member obtains. The descent is not a sequence of conditions waiting to be satisfied, in which the first member would obtain only once the later ones had been secured; it is a sequence of facts, every one of which obtains, and each of which is accounted for by the next. Nothing is postponed, because the relation is not temporal and is not a relation of conditional supply; it is a relation of metaphysical explanation, and every member has its explanation. The temptation to think otherwise comes from an image of grounding as the transmission of a quantity, \enquote{reality} or \enquote{being}, from ground to grounded, which we examine in Objection~\ref{o:borrowed}.

\ref{p:stock} helps to see why the descent is harmless. In the descent, nothing is borrowed in the sense in which a derivative property borrows its nature from the properties it is constructed from. The mass of a cube at level $n$ is of the same kind as the masses of its sub-cubes; the descent never asks where mass comes from, only how much of it is where. A descent of \emph{properties}, in which every property was constructed from more basic properties without end, would be a different matter, since every property would then be characterized only by reference to others and none by what it is; we do not assert such a descent, and \ref{p:stock} denies it.

\section{The descent and the cosmos}\label{s:cosmos}

MET-09 argues that the cosmos is the only fundamental object and that ordinary objects are abstractions from it. On its face, this conflicts with the present thesis: one manuscript says that something is fundamental, the other that nothing is. The conflict is only apparent, and it is worth making the reconciliation explicit, since the two theses together form a single picture.

The present thesis concerns levels, and a level, by Definition~\ref{d:level}, consists of proper parts of the cosmos. The descent concerns proper parts: there is no smallest level of parts, and every level is grounded in a finer one. The thesis of MET-09 concerns the whole: the fact that the cosmos exists and is as it is, which is not a fact at any level, is not grounded in anything, and every fact about a proper part is partly grounded in it. Nothing in Theorem~\ref{t:nofund} says that facts at a level are grounded \emph{only} in facts at finer levels, and nothing in the thesis of MET-09 says that the facts about proper parts are not also grounded in facts about their own parts. The two theses therefore assert different groundings, and the question is only whether they can hold together.

\begin{proposition}[Compatibility]\label{p:compat}
There is a strict partial order on facts in which (i) every level is grounded in a lower one and no level is fundamental, (ii) there is an infinite descending sequence of facts at successively finer levels, and (iii) there is exactly one fact not partly grounded in any fact, a fact about the cosmos, and it partly grounds every fact at every level.
\end{proposition}

\begin{proof}
Extend the model of Proposition~\ref{p:stockcons} with one further fact $c$, that the cosmos exists and is as it is, and let $c$ partly ground every fact at every level, and nothing partly ground $c$. The resulting relation is irreflexive, since $c$ does not ground itself and the original relation is irreflexive. It is transitive: a chain $c \prec f \prec g$ gives $c \prec g$, which holds since $c$ grounds every level fact, and chains among level facts are as before. (i) and (ii) hold as in the original model, since full grounds at finer levels are unaffected. (iii) holds by construction: every level fact is partly grounded in $c$, and $c$ in nothing.
\end{proof}

In the order of Proposition~\ref{p:compat}, the cosmos is the largest object and the least element of the grounding order. Mereologically, the levels descend from the cosmos toward ever smaller parts with no smallest; in the order of grounding, they descend toward the cosmos, which bounds them all from below without being one of them. Grounding is not well-founded in the sense that excludes infinite descending chains; but it meets a weaker condition, that every non-fundamental fact is partly grounded in something fundamental, of the kind that \citet{Dixon2016} shows must be distinguished from the absence of such chains. Priority monism, so combined with infinitism about levels, is a version of the combination that \citet{Schaffer2010} defended: he held that gunk is possible and that the whole can be fundamental when no part is. Our version differs from his in one respect. Schaffer took the possibility of gunk to be an argument for monism, on the ground that there must be a fundamental basis and that in a gunky world only the whole can supply one. We have argued that the strong requirement is not needed: the descent is not vicious whether or not the cosmos is fundamental, as Proposition~\ref{p:stockcons}, whose model contains no fundamental fact, shows. Whether the cosmos is in fact fundamental is a further question, which MET-09 answers on other grounds.

\section{Further cases}\label{s:cases}

\paragraph{Causal drainage.} \citet{Block2003} argued that if the causal powers of each level are excluded by those of the level beneath, and there is no bottom level, there are no causal powers anywhere. The argument requires the exclusion principle in a form on which a sufficient cause at a lower level excludes any distinct cause at a higher one. If the exclusion principle is true in that form, it drains causal powers from every level above the bottom whether or not there is a bottom, and so refutes itself against the evident efficacy of tables and organisms; if it is false, there is no drainage. The descent of levels plays no role. On the account of causation of MET-04, a cause is a counterfactual difference-maker selected by norms of normality, and difference-makers are found at many levels at once: the push of the hand and the motion of the molecules of the hand both make a difference to the motion of the table. Exclusion in Block's form is false on that account, and the descent drains nothing.

\paragraph{Mixed hierarchies.} Some regions of the actual world may be structured differently from others. Suppose, contrary to the argument of Section~\ref{s:divisibility}, that some objects are simples. Then a level may contain simples, and a finer level is still available as long as some member of every level has proper parts; \ref{p:divisibility} requires only that, and the thesis survives. What would refute it is a level consisting entirely of simples, a level of atoms in the sense of the atomists, and that is what the continuity of fields excludes.

\paragraph{Ascending regresses.} Not every infinite regress is a descent of grounds. The truth regress, from $p$ to \enquote{it is true that $p$} and onwards, ascends: each member is grounded in the one before it, and the first is grounded in whatever grounds $p$. The instantiation regress of MET-22-a is of the same kind. Ascending regresses are benign for a reason that does not apply to descents, namely that every member has a finite grounding history; our argument shows that this reason, though sufficient, is not necessary.

\paragraph{Justification.} EPI-05 argues that justification has a tree structure whose roots are perceptual experiences requiring no further support, and so that epistemic justification is well-founded. That is no tension. Justification is a relation among a subject's states, which are finite in number; grounding among levels is a relation among facts about the world, which need not be. An argument that a regress of reasons must terminate in a finite mind is not an argument that a regress of grounds must terminate in the world.

\section{Objections and replies}\label{s:objections}

\begin{objection}[Borrowed being]\label{o:borrowed}
In a descending sequence, each fact has its reality only derivatively, from the next. If every fact has its reality only derivatively, reality is borrowed at every step and never possessed by anything in its own right; nothing then has any reality to lend. The descent is like a chain of loans with no lender \citep[cf.][]{Schaffer2010}.
\end{objection}

\begin{reply}
The analogy presupposes that grounding transmits a quantity, as a loan transmits money, so that the grounded has what it has only because the ground had it first. Grounding is not like that. A conjunction obtains in virtue of its conjuncts, but the conjuncts do not lend the conjunction a portion of their obtaining; each obtains, and the conjunction obtains because they do. The same holds of a table and its parts. Each fact in the descent obtains, and each has a full ground in facts that obtain. The thought that some fact must obtain \enquote{in its own right} is the thought that some fact must be ungrounded, which is the conclusion at issue; it is not a premise that can be used to establish it. If the image of borrowing is taken literally, it equally condemns the foundationalist's derivative facts, which also have their reality \enquote{only derivatively}; if it is not taken literally, it has no force.
\end{reply}

\begin{objection}[Unified explanation]\label{o:unified}
Even if each fact in the descent is explained, the totality of derivative facts is not. The foundationalist can say why there is anything derivative at all: because of the fundamental facts. The infinitist cannot. So foundationalism has an explanatory advantage, which is a reason to believe it \citep{Cameron2008}.
\end{objection}

\begin{reply}
The foundationalist explains the totality of derivative facts by the fundamental facts, and leaves the fundamental facts unexplained. The infinitist explains every fact by other facts and leaves the totality of level facts unexplained. Neither explains everything, and the question is which leaves less unexplained. On the foundationalist picture, every fact at the bottom level is a brute fact; on the infinitist picture, no fact at any level is. What is left unexplained on the infinitist picture is only the existence of the whole system of levels, and that, on the foundationalist picture, is left unexplained as well, since it consists in the fundamental facts and what they ground. The explanatory advantage, if any, therefore lies with the infinitist. If, in addition, MET-09 is right that the cosmos is fundamental, the totality of level facts is partly grounded in a single fact about the cosmos, and the unified explanation the objection asks for is supplied, without a bottom level.
\end{reply}

\begin{objection}[Fundamental properties]\label{o:props}
The sparse theory of MET-02 says that all properties are constructed from genuine ones. Genuine properties are fundamental properties, and fundamental properties are properties of the fundamental level; so the sparse theory requires a fundamental level.
\end{objection}

\begin{reply}
The objection equivocates between two conceptions of sparse properties \citep{Schaffer2004}. On the fundamental conception, the genuine properties are by definition those of the bottom level; on the scientific conception, they are the properties that ground objective resemblance, wherever they are instantiated. MET-02-a adopts the scientific conception, and on it the genuine properties can be instantiated at every level, as \ref{p:stock} and Proposition~\ref{p:stockcons} show. The descent concerns objects, not properties.
\end{reply}

\begin{objection}[Strong emergence]\label{o:emergence}
Some wholes may have features not grounded in their parts and their relations: consciousness, on some views, or the states of entangled systems. Then \ref{p:aggregation} fails, and some facts at higher levels are not grounded in lower ones.
\end{objection}

\begin{reply}
The objection, if successful, would refute the claim that every level is grounded in a lower one, but it would not restore a fundamental level: it would make some facts at some levels ungrounded by finer levels, while every level still has finer levels beneath it. In any case we do not accept the premise. Entangled states are features of the whole system, but they are captured in the facts at a finer level by the non-separable relations among the parts, which Definition~\ref{d:facts} includes; and the states of parts and their relations are, if MET-09 is right, themselves partly grounded in the state of the cosmos. Consciousness, on the view of MND-01, is the intrinsic nature of physical structure, not an addition to it; whatever grounds the structure at one level grounds its nature as well.
\end{reply}

\section{What the result does not establish}\label{s:limits}

The result does not establish that infinite descent is necessary: a world with a fundamental level is possible, and the argument for \ref{p:divisibility} is a posteriori. It does not establish that nothing is fundamental. It is consistent with the thesis of MET-09 that the cosmos is fundamental, as Proposition~\ref{p:compat} shows, and with the existence of fundamental facts that are not about any level. It does not establish a descent of properties; it requires, on the contrary, that the stock of genuine properties be fixed. It does not establish that grounding is a single relation, which is the thesis of MET-14, although it uses that thesis to read the transitivity of grounding without restriction. And it does not show that the special sciences cannot treat a level as fundamental for their purposes; the sense in which fundamentality may be relative to an inquiry is the subject of MET-10-b.

\section{Conclusion}

There is no fundamental level of reality. The facts about the objects at any level are grounded in the facts about their parts and the relations among them, and the continuity of the physical world ensures that every level of parts has a finer level. The resulting infinite descent is consistent, leaves no fact without a full ground, and does not answer any general question by presupposing its answer; it is therefore not vicious. The descent concerns objects, not properties, and it is compatible both with a fixed stock of genuine properties at every level and with the fundamentality of the cosmos as a whole.

\begin{thebibliography}{99}\small
\bibitem[Bliss(2013)]{Bliss2013} Bliss, R.~L. (2013). Viciousness and the structure of reality. \emph{Philosophical Studies}, 166(2), 399--418.
\bibitem[Block(2003)]{Block2003} Block, N. (2003). Do causal powers drain away? \emph{Philosophy and Phenomenological Research}, 67(1), 133--150.
\bibitem[Butterfield(2006)]{Butterfield2006} Butterfield, J. (2006). Against \emph{pointillisme} about mechanics. \emph{The British Journal for the Philosophy of Science}, 57(4), 709--753.
\bibitem[Cameron(2008)]{Cameron2008} Cameron, R.~P. (2008). Turtles all the way down: Regress, priority and fundamentality. \emph{The Philosophical Quarterly}, 58(230), 1--14.
\bibitem[Dixon(2016)]{Dixon2016} Dixon, T.~S. (2016). What is the well-foundedness of grounding? \emph{Mind}, 125(498), 439--468.
\bibitem[Fine(2012)]{Fine2012} Fine, K. (2012). Guide to ground. In F.~Correia and B.~Schnieder (eds.), \emph{Metaphysical Grounding: Understanding the Structure of Reality}. Cambridge University Press.
\bibitem[Kim(1998)]{Kim1998} Kim, J. (1998). \emph{Mind in a Physical World}. MIT Press.
\bibitem[Lewis(1991)]{Lewis1991} Lewis, D. (1991). \emph{Parts of Classes}. Blackwell.
\bibitem[Morganti(2014)]{Morganti2014} Morganti, M. (2014). Metaphysical infinitism and the regress of being. \emph{Metaphilosophy}, 45(2), 232--244.
\bibitem[Nolan(2001)]{Nolan2001} Nolan, D. (2001). What's wrong with infinite regresses? \emph{Metaphilosophy}, 32(5), 523--538.
\bibitem[Oppenheim and Putnam(1958)]{OppenheimPutnam1958} Oppenheim, P. and Putnam, H. (1958). Unity of science as a working hypothesis. In H.~Feigl, M.~Scriven and G.~Maxwell (eds.), \emph{Minnesota Studies in the Philosophy of Science}, Vol.~2. University of Minnesota Press.
\bibitem[Rosen(2010)]{Rosen2010} Rosen, G. (2010). Metaphysical dependence: Grounding and reduction. In B.~Hale and A.~Hoffmann (eds.), \emph{Modality: Metaphysics, Logic, and Epistemology}. Oxford University Press.
\bibitem[Russell(1918)]{Russell1918} Russell, B. (1918--19). The philosophy of logical atomism. \emph{The Monist}, 28--29.
\bibitem[Schaffer(2003)]{Schaffer2003} Schaffer, J. (2003). Is there a fundamental level? \emph{No\^us}, 37(3), 498--517.
\bibitem[Schaffer(2004)]{Schaffer2004} Schaffer, J. (2004). Two conceptions of sparse properties. \emph{Pacific Philosophical Quarterly}, 85(1), 92--102.
\bibitem[Schaffer(2009)]{Schaffer2009} Schaffer, J. (2009). On what grounds what. In D.~Chalmers, D.~Manley and R.~Wasserman (eds.), \emph{Metametaphysics}. Oxford University Press.
\bibitem[Schaffer(2010)]{Schaffer2010} Schaffer, J. (2010). Monism: The priority of the whole. \emph{The Philosophical Review}, 119(1), 31--76.
\bibitem[Streater and Wightman(1964)]{StreaterWightman1964} Streater, R.~F. and Wightman, A.~S. (1964). \emph{PCT, Spin and Statistics, and All That}. W.~A. Benjamin.
\bibitem[Tahko(2014)]{Tahko2014} Tahko, T.~E. (2014). Boring infinite descent. \emph{Metaphilosophy}, 45(2), 257--269.
\bibitem[Zimmerman(1996)]{Zimmerman1996} Zimmerman, D.~W. (1996). Could extended objects be made out of simple parts? An argument for \enquote{atomless gunk}. \emph{Philosophy and Phenomenological Research}, 56(1), 1--29.
\end{thebibliography}

\end{document}
